




\section{Setup and statement}
Let $Q=\{1,\ldots,d\}$, and let $R$ and $B$ be transformation matrices
on $Q$: every row contains exactly one entry equal to $1$, all other
entries being zero. We use right-action composition,
\[
 x(fg)=(xf)g.
\]
Thus a transformation $f$ has matrix entry $f_{xy}=1$ exactly when
$xf=y$. Its matrix rank equals $|\operatorname{im} f|$.

Fix $0<p<1$ and set
\[
 T=pR+(1-p)B,\qquad S=\langle R,B\rangle.
\]
Assume $T$ is irreducible. Let $\pi$ be its stationary probability row
vector, so that
\[
 \pi T=\pi,\qquad \pi\mathbf{1}=1.
\]
Write
\[
 r=\min_{s\in S}\rank(s),\qquad K=\{s\in S:\rank(s)=r\}.
\]
The set $K$ is the minimal ideal (kernel) of this finite transformation
semigroup. In particular, the word ``kernel'' here denotes a semigroup
ideal, whereas a transformation's fibers form its kernel partition.

Define the linear operator
\[
 \Phi(X)=pRXR^{\mathsf{T}}+(1-p)BXB^{\mathsf{T}}
\]
on real $d\times d$ matrices. Identifying $X$ with its array of entries
indexed by ordered pairs, this is the operator with transition matrix
\[
 P_2=p(R\otimes R)+(1-p)(B\otimes B).
\]
Both particles use the same randomly chosen letter; this is not the
independent-particle transition matrix $T\otimes T$.

\begin{theorem}\label{tensor-book:thm:main}
% BEGIN SUBJECT INDEX: tensor-book:thm:main
\index{equipartition!tensor-square proof}
% END SUBJECT INDEX

The Ces\`aro limit
\[
 H=\lim_{N\to\infty}\frac1N\sum_{\ell=0}^{N-1}\Phi^\ell(I_d)
\]
exists and satisfies
\[
 H\succeq0,\qquad H_{xx}=1,\qquad \Phi(H)=H.
\]
Moreover,
\begin{equation}\label{tensor-book:eq:stationary}
 H\pi^{\mathsf{T}}=\frac1r\mathbf{1},
 \qquad
 \pi H\pi^{\mathsf{T}}=\frac1r.
\end{equation}
For every $k\in K$ and every fiber $F$ of $k$,
\begin{equation}\label{tensor-book:eq:friedman}
 \pi(F):=\sum_{x\in F}\pi_x=\frac1r.
\end{equation}
Consequently the kernel rank is determined by the tensor-square action:
\[
 \boxed{\displaystyle r=\frac{1}{\pi H\pi^{\mathsf{T}}}.}
\]
\end{theorem}

Equation~\eqref{tensor-book:eq:friedman} is the equal-stationary-mass conclusion
associated with Friedman's work on road coloring~\cite{F}.
The proof below is self-contained.

\section{The collision matrix}
For a word $w$ in the two letters, write $W$ for its transformation
matrix, and let $q_w$ be the product of its letter probabilities. Thus
$q_w>0$. Summing over words (not merely over distinct resulting matrices),
\begin{equation}\label{tensor-book:eq:word-average}
 \Phi^\ell(X)=\sum_{|w|=\ell}q_wWXW^{\mathsf{T}},
 \qquad
 T^\ell=\sum_{|w|=\ell}q_wW.
\end{equation}
At length zero, the word matrix is $I_d$.

\begin{lemma}\label{tensor-book:lem:H}
The matrix $H$ in Theorem~\ref{tensor-book:thm:main} exists, is positive
semidefinite, has diagonal entries one, and satisfies $\Phi(H)=H$.
If distinct states $x,y$ cannot be merged by any word, then $H_{xy}=0$.
\end{lemma}
\begin{proof}
The matrix $P_2$ is a finite stochastic matrix. Its Ces\`aro powers
converge, so their action on the vectorized identity also converges.
This proves existence of $H$.

Each $WW^{\mathsf{T}}$ is positive semidefinite, and
\[
 (WW^{\mathsf{T}})_{xy}=\begin{cases}1,&xW=yW,\\0,&xW\ne yW.\end{cases}
\]
In particular its diagonal entries equal one. Formula
\eqref{tensor-book:eq:word-average} and passage to the limit prove the corresponding
assertions for $H$. The fixed-point equation follows either from the
Ces\`aro projection or from the telescoping identity
\[
 \Phi\left(\frac1N\sum_{\ell=0}^{N-1}\Phi^\ell(I_d)\right)
 -\frac1N\sum_{\ell=0}^{N-1}\Phi^\ell(I_d)
 =\frac{\Phi^N(I_d)-I_d}{N}.
\]
The right side tends to zero because all its entries are bounded in
absolute value by $1/N$. Finally, a nonmergeable pair has zero entry in
every matrix being averaged, and hence in $H$.
\end{proof}

\begin{remark}[Probabilistic interpretation]
The entry $(\Phi^\ell(I_d))_{xy}$ is the probability that two particles
started at $x,y$ agree after $\ell$ common random letters. Once they
agree, they remain together. These probabilities are therefore
nondecreasing in $\ell$ and bounded by one. Thus, in this particular
construction, the ordinary limit $\Phi^\ell(I_d)\to H$ also exists,
and $H_{xy}$ is the probability of eventual coalescence. Ces\`aro
averaging provides a direct finite-Markov-chain construction that does
not require aperiodicity.
\end{remark}

\section{Positive semidefiniteness and stationarity}
\begin{lemma}\label{tensor-book:lem:word-invariance}
For every word matrix $W\in S$,
\begin{equation}\label{tensor-book:eq:word-invariance}
 HW^{\mathsf{T}}\pi^{\mathsf{T}}=H\pi^{\mathsf{T}}.
\end{equation}
Furthermore, if $c=\pi H\pi^{\mathsf{T}}$, then
\begin{equation}\label{tensor-book:eq:constant}
 H\pi^{\mathsf{T}}=c\mathbf{1}.
\end{equation}
\end{lemma}
\begin{proof}
Fix a word length $\ell\ge1$. Since $\Phi^\ell(H)=H$ and
$\pi T^\ell=\pi$, the word averages satisfy
\[
 \sum_{|w|=\ell}q_w\pi W=\pi,
 \qquad
 \sum_{|w|=\ell}q_w(\pi W)H(\pi W)^{\mathsf{T}}=c.
\]
Expanding the squared deviations therefore gives
\begin{equation}\label{tensor-book:eq:variance}
 \sum_{|w|=\ell}q_w
 (\pi W-\pi)H(\pi W-\pi)^{\mathsf{T}}=0.
\end{equation}
Every summand is nonnegative because $H\succeq0$, and every $q_w$ is
positive. For a positive semidefinite matrix, $vHv^{\mathsf{T}}=0$ implies
$Hv^{\mathsf{T}}=0$. Hence
\[
 H(\pi W-\pi)^{\mathsf{T}}=0
\]
for every word of length $\ell$. As $\ell$ was arbitrary, this proves
\eqref{tensor-book:eq:word-invariance} for every word.

Set $h=H\pi^{\mathsf{T}}$. Applying the fixed-point equation and the
length-one cases of \eqref{tensor-book:eq:word-invariance}, we obtain
\[
 \begin{aligned}
 h&=pRH R^{\mathsf{T}}\pi^{\mathsf{T}}+(1-p)BH B^{\mathsf{T}}\pi^{\mathsf{T}}\\
  &=pRh+(1-p)Bh=Th.
 \end{aligned}
\]
For an irreducible finite stochastic matrix, the right fixed vectors
are the constant vectors. Thus $h=\alpha\mathbf{1}$. Multiplying by $\pi$
gives $\alpha=\pi h=c$, proving \eqref{tensor-book:eq:constant}.
\end{proof}

\begin{remark}
To obtain \eqref{tensor-book:eq:word-invariance} for arbitrary words, the variance
identity is applied separately at every word length. The length-one
identities alone do not justify a direct iteration without a further
invariance argument.
\end{remark}

\section{Minimum-rank ranges and equal fiber masses}
\begin{lemma}\label{tensor-book:lem:range}
For every $k\in K$,
\begin{equation}\label{tensor-book:eq:range}
 kHk^{\mathsf{T}}=kk^{\mathsf{T}}.
\end{equation}
\end{lemma}
\begin{proof}
Let $U=\operatorname{im} k$, so $|U|=r$. No word can merge two distinct points of
$U$: otherwise its composition after $k$ would have rank below $r$.
Lemma~\ref{tensor-book:lem:H} consequently gives
\[
 H_{uv}=\delta_{uv}\qquad(u,v\in U).
\]
For any $x,y\in Q$, therefore,
\[
 (kHk^{\mathsf{T}})_{xy}=H_{xk,yk}
 =\begin{cases}1,&xk=yk,\\0,&xk\ne yk,\end{cases}
 =(kk^{\mathsf{T}})_{xy}.
\]
\end{proof}

\begin{proof}[Proof of Theorem~\ref{tensor-book:thm:main}]
The initial properties of $H$ were proved in Lemma~\ref{tensor-book:lem:H}.
For $k\in K$, Lemmas~\ref{tensor-book:lem:word-invariance} and~\ref{tensor-book:lem:range}
give
\[
 \begin{aligned}
 kk^{\mathsf{T}}\pi^{\mathsf{T}}
 &=kHk^{\mathsf{T}}\pi^{\mathsf{T}}\\
 &=kH\pi^{\mathsf{T}}\\
 &=c\,k\mathbf{1}=c\mathbf{1},
 \end{aligned}
\]
where $c=\pi H\pi^{\mathsf{T}}$. On the other hand,
\[
 (kk^{\mathsf{T}}\pi^{\mathsf{T}})_x
 =\sum_{y:\,yk=xk}\pi_y.
\]
This is precisely the stationary mass of the fiber containing $x$.
Thus all fibers of $k$ have mass $c$. There are $r$ such fibers and
they partition $Q$, so $rc=1$. This proves all the remaining assertions.
\end{proof}

\section{Consequences and scope}
\begin{corollary}[A coloring-independent arithmetic restriction]
Suppose $\pi_x=w_x/W$, where the $w_x$ are positive integers and
$W=\sum_xw_x$. Then every kernel block has integer weight $W/r$;
in particular $r$ divides $W$.
If $\pi$ is uniform, every kernel block has cardinality $d/r$, so
$r$ divides $d$.
\end{corollary}
\begin{proof}
Multiply the equal-mass identity by $W$, or by $d$ in the uniform case.
\end{proof}

\begin{corollary}[Synchronization]
The following are equivalent:
\begin{enumerate}
\item $S$ contains a constant transformation;
\item $r=1$;
\item $H=\mathbf{1}\mathbf{1}^{\mathsf{T}}$;
\item $\pi H\pi^{\mathsf{T}}=1$.
\end{enumerate}
\end{corollary}
\begin{proof}
The equivalence of the first two statements follows from the definition
of rank, and their equivalence with the fourth follows from the rank
formula. A reset word occurs eventually with probability one among
independent random letters, for instance by considering disjoint blocks
of its length. It therefore forces eventual coalescence of every pair,
so the first statement implies the third. The third implies the fourth
because $\pi\mathbf{1}=1$.
\end{proof}

\begin{remark}[Application to road colorings]
For a $0$--$1$ adjacency matrix $A$ having two ones per row and a coloring
$A=R+B$, take $p=1/2$. Then $T=A/2$ and its stationary distribution do
not depend on the coloring, whereas $P_2$ and $H$ generally do.
Thus every coloring obeys the same stationary-weight restrictions,
while its particular tensor-square fixed point determines its kernel
rank. Irreducibility and aperiodicity are the standing assumptions for
the road-coloring existence problem; only irreducibility is used in the
equal-mass theorem proved here.
\end{remark}

\begin{remark}[More than two generators]
The entire proof applies to finitely many transformation matrices
$M_1,\ldots,M_s$ with positive probabilities $p_i$ summing to one:
replace $T$ by $\sum_i p_iM_i$ and $\Phi(X)$ by
$\sum_i p_iM_iXM_i^{\mathsf{T}}$. The same conclusions hold whenever $T$
is irreducible.
\end{remark}


